% !TEX program = xelatex
\documentclass[tikz,border=8pt]{standalone}
\usepackage[UTF8]{ctex}
\usetikzlibrary{arrows.meta,positioning,calc,shadows.blur}
\definecolor{blueA}{HTML}{2454A6}
\definecolor{purpleA}{HTML}{7B61B8}
\definecolor{redA}{HTML}{D95D5D}
\definecolor{darkA}{HTML}{243043}
\definecolor{grayA}{HTML}{6D7788}
\definecolor{lightA}{HTML}{F4F7FB}
\definecolor{noteA}{HTML}{FFF8EC}
\definecolor{noteB}{HTML}{F1D6A5}
\begin{document}
\begin{tikzpicture}[font=\sffamily]
\node[font=\bfseries\Large, text=darkA] at (7.5,5.55) {低频与高频信息的视觉语义示意};
\node[font=\small, text=grayA] at (7.5,5.1) {低频偏向全局结构，高频偏向边缘、纹理与噪声};
\tikzset{panel/.style={fill=lightA, rounded corners=9pt, minimum width=3.8cm, minimum height=3.2cm, blur shadow}}
\node[panel] (p1) at (2.7,2.95) {};
\node[panel] (p2) at (7.5,2.95) {};
\node[panel] (p3) at (12.3,2.95) {};
\node[font=\bfseries, text=darkA] at (2.7,4.3) {低频分量};
\node[font=\bfseries, text=darkA] at (7.5,4.3) {高频分量};
\node[font=\bfseries, text=darkA] at (12.3,4.3) {混合视觉信号};
\draw[blueA, line width=2pt, smooth, domain=0:3.0, samples=120, variable=\x] plot ({1.2+\x}, {2.95+0.55*sin(240*\x)});
\draw[redA, line width=1.8pt, smooth, domain=0:3.0, samples=180, variable=\x] plot ({6.0+\x}, {2.95+0.30*sin(1800*\x)+0.17*sin(2800*\x)});
\draw[purpleA, line width=1.8pt, smooth, domain=0:3.0, samples=180, variable=\x] plot ({10.8+\x}, {2.95+0.48*sin(240*\x)+0.18*sin(1800*\x)});
\node[font=\small, text=grayA] at (2.7,1.55) {轮廓 · 光照 · 全局布局};
\node[font=\small, text=grayA] at (7.5,1.55) {边缘 · 纹理 · 噪声细节};
\node[font=\small, text=grayA] at (12.3,1.55) {结构与细节共同形成感知};
\node[fill=noteA, draw=noteB, rounded corners=7pt, minimum width=12.8cm, minimum height=0.7cm, text=darkA, font=\small] at (7.5,0.55) {综述写作提示：分析方法时应说明其关注的频带、作用位置、带来的性能收益及可能损失的信息。};
\end{tikzpicture}
\end{document}
